\begin{solucion}
(b) Para \(x^{4}-2\).  
De manera análoga, \(x^{4}-2\) es irreducible sobre \(\mathbb{Q}\) (por el criterio de Eisenstein con \(p=2\), es irreducible en \(\mathbb{Q}[x]\)).  
Si tomamos \(\beta=\sqrt[4]{2}\) (raíz real de \(x^{4}-2=0\)), obtenemos primero la extensión \(\mathbb{Q}(\beta)\) de grado \(4\).  
En \(\mathbb{Q}(\beta)\) el polinomio se factoriza como \((x-\beta)(x+\beta)\), faltando incorporar las raíces complejas \(\pm i\beta\).  
Al igual que en el inciso anterior, \(i\notin\mathbb{Q}(\beta)\), así que debemos adjuntarlo.  
El cuerpo de descomposición buscado es
\[
M=\mathbb{Q}(\beta,i)=\mathbb{Q}\!\bigl(\sqrt[4]{2},\,i\bigr),
\]
Además,
\[
[\mathbb{Q}(\beta,i):\mathbb{Q}]
  =[\mathbb{Q}(\beta,i):\mathbb{Q}(\beta)]
  \cdot[\mathbb{Q}(\beta):\mathbb{Q}]
  =2\cdot4=8.
\]

Por lo tanto,
\[
\text{el cuerpo de descomposición de }x^{4}-2\text{ es }\mathbb{Q}\bigl(\sqrt[4]{2},\,i\bigr)
\text{ y su grado sobre }\mathbb{Q}\text{ es }8.
\]
\end{solucion}
